\HMTNarrativeHeading{De la estructura determinantal a la orientación}

Las constantes de Planck se presentan después de los objetos que
determinan sus secciones de acción. El registro dodecafásico contiene el
bloque local de Hadamard; su cociente integral tiene un orden calculable.
El lector determinantal compone ese dato con las coordenadas ya generadas
y fija una valoración decimal. La escala y el coeficiente de una sección
cumplen funciones diferentes, que la demostración mantiene explícitas.

El retorno distingue las secciones anterior y posterior a la vuelta.
Sus coeficientes conservan la procedencia del carácter y de la
continuación regional. La expresión de la acción en unidades físicas
se introduce mediante la carta dimensional correspondiente. La
construcción transmite así una sección, su coeficiente, su escala y la
regla de lectura. El uso posterior de la constante conserva esas
relaciones.

Ángulo y contraángulo se desarrollan a continuación como coordenadas de
orientación. La primera lectura desplaza los bloques de alfa en la base
de completación; la segunda compone el coeficiente adimensional de
retorno con alfa. La carta circular expresa el resultado en grados o
radianes y conserva el número de vueltas cuando se utiliza el
levantamiento. Esta secuencia explica por qué las coordenadas angulares
son intermediarias importantes entre la acción y los lectores
geométricos posteriores.

Los canales conjugados reciben esa orientación y evalúan, junto con
la incidencia marcada, el funcional de Barbero--Immirzi.
La realización electrónica utiliza su trayectoria y sus lectores de
retorno, y recibe el cociente de las secciones de acción. La narración conserva
el orden de sus antecedentes y distingue el coeficiente adimensional
de la magnitud de acción que contribuye a representar. Una misma
constante puede actuar en varias etapas sin perder el mapa por el que
fue obtenida.

El recorrido prepara una composición física concreta. Acción y reloj
determinan una energía; la transducción térmica relaciona esa energía
con la información; el lector de área incorpora la incidencia marcada.
La realización gravitatoria completa después el cociente entre energía
y área. El artículo conserva los cálculos de esas composiciones para
mostrar cómo los resultados aritméticos y geométricos participan en
una lectura física común.
